\documentclass{article}
\usepackage{xspace}
\newcommand{\latex}{\LaTeX\xspace}

\usepackage{graphicx} % need to include the 'graphicx' package for figures.
\graphicspath{{./Figs/}} % Tell LaTeX where it should look for your figures.

% For groups of figure folders you can do this...
% \graphicspath{{subdir1/}{subdir2/}{subdir3/}...{subdirn/}} 


\author{Dylan Mikesell} % define the author
\title{\latex Tutorial} % define the title
\date{01 October 2015} % set the date (for today you can use the command \today instead of explicitly writing the date.)

% end the PREAMBLE
%------------------------------------------------------------------------------
\begin{document}

\maketitle % generate the title

\newpage

\tableofcontents % make a table of contents

\newpage

\section{My first section} % make a section for the table of contents

This is my first \latex document. And this is my first figure reference (Fig.~\ref{fig:firstFigure}).% type some words

\begin{figure}
	\centering
	\includegraphics[width=0.5\columnwidth]{LaTeXLogo.png}
	% \includegraphics[width=0.5\columnwidth]{./Figs/LaTeXLogo.png}
	\caption{This is my first figure.}
	\label{fig:firstFigure}
\end{figure}

This is my first table (Tab.~\ref{tab:firstTable}).

% what does the \noindent command do? Uncomment line below.
% \noindent This is my first table (Tab.~\ref{tab:firstTable}).

\begin{table}[b]
	% \centering
	\begin{tabular}{l|c|r}
	$a \approx b$ & $a \ne b$ & $a = \int_{t_0}^{t_1} \frac{g(t)}{t} dt $ \\ \hline
	\textbf{column 1} & \textbf{column 2} & \textbf{column 3} \\ \hline
	\end{tabular}
	\caption{I just made my first table!}
	\label{tab:firstTable}
\end{table}

\newpage

\section{Adios} % make a final section

\ldots{} and that's it.

\end{document}